% DSRM phases and objectives

\section{Method}
\label{sec:method}

The work follows Design Science Research Methodology \cite{peffers2007dsrm}. The main output is a running platform, which is an instantiation artifact in the design science taxonomy \cite{hevner2004design}, so an artifact-oriented framework fits better than a pure software engineering process. CRISP-DM and TDSP were considered and set aside, because both organise the building of one model on one data domain, whereas the output here is a cross-domain platform on which a model is only one of the workloads served.

The six phases run iteratively. Problem identification produces four problem statements from a study of DataOps, MLOps, and data governance. Objective definition turns each statement into an artifact that can be installed, operated, and audited, giving four objectives mapped one to one onto the problem statements.

\begin{itemize}
\item \textbf{T-1} unifies ingestion, processing, storage, feature store, training, serving, governance, and observability into one declarative control plane with GitOps practice.
\item \textbf{T-2} builds a data governance sub-system providing metadata catalog, lineage recording, and quality control as shared services along the pipeline.
\item \textbf{T-3} implements automated drift detection with a threshold-based retraining policy, and enforces point-in-time correctness across batch and streaming modes through the historical feature join.
\item \textbf{T-4} develops a dual-store feature service delivering tabular and aggregate features under one API contract, with vector embeddings served on a separate index whose embedding space stays consistent between training and inference.
\end{itemize}

The objectives are decomposed into 14 functional and 12 non-functional requirements, which become the units of evaluation in Section~\ref{sec:evaluation}. Design and development classify the services named by the sources in Section~\ref{subsec:layers} into platform layers, select open source components per layer, and express the result as a declarative control plane. Before that, four provisioning alternatives were scored against the same 26 requirements from vendor documentation and public offerings, and open source on self-managed Kubernetes scored highest, mainly on temporal validity, vector support, portability, and extensibility. That scoring is desk research preceding implementation and is not an evaluation result. It is also separate from the three-way cost comparison in Section~\ref{subsec:cost}, which prices only the options that remained practical. Demonstration installs the platform on k3s and drives it with one real domain use case. Evaluation examines functional fulfilment, non-functional targets, and governance, and communication produces this paper alongside the source repository.
